% Generated. Do not edit.
% Projection scriptural-date-location/section, version 2: the Scriptural Date and
% Location page, generated whole from research/chronology.toml, the Scripture
% chronology corpus and the source library by tools/tpt document-types refresh.
% document: liturgy/roman-rite/postconciliar/roman-missal-third-edition-en-us-2011/propers/temporal/pc-s01-first-sunday-of-advent-year-b
\runninghead{Scriptural Date and Location}
\section{Scriptural Date and Location}
\label{sec:date-location}

{\small Each passage the Mass reads or sings from Scripture, in canonical order. The Context column distinguishes the event or setting, an attributed author's era, and the writing or formation of the text. The Date column gives calendar orientation. A proposal span summarizes the extent of numerical dates in the sources, including gaps between proposals. Undated occasions and events remain distinct from composition dates. The complete alternatives and their evidence follow the table.\par}
\medskip
% Generated. Do not edit.
% Deterministic projection of research/chronology.toml.
% schema: 5
% document: liturgy/roman-rite/postconciliar/roman-missal-third-edition-en-us-2011/propers/temporal/pc-s01-first-sunday-of-advent-year-b
% generated-by: tools/tpt proper-chronology annotations
\newcommand{\chronologyannotationclaim}[6]{#6}
\newcommand{\chronologyannotationcomparisonclaim}[8]{#8}
\newcommand{\chronologyannotationreach}[3]{}
\newcommand{\chronologyannotationgroup}[3]{#3}
\newcommand{\chronologyannotationcomparisongroup}[4]{#4}
\newcommand{\chronologyannotation}[1]{%
  \ifcsname triptychchronologyannotation@#1\endcsname
    \csname triptychchronologyannotation@#1\endcsname
  \else
    \PackageError{triptych}{No chronology annotation for #1}{}%
  \fi}
% appointment-dependency: guidance/liturgy/postconciliar-propers-registry.md sha256=376860a7c37a4c75bb6d880cd1b260183be1877a07e3cd01660a9e2854e2ea72
% appointment-dependency: src/gpt/liturgy/roman-rite/postconciliar/roman-missal-third-edition-en-us-2011/propers/registry/README.md sha256=b11e73fa4ce7ca027f7ffd13cc013c843c1f266bdbd35aa4958f4ff494f002b4
% appointment-dependency: src/gpt/liturgy/roman-rite/postconciliar/roman-missal-third-edition-en-us-2011/propers/registry/formula-dispositions.md sha256=efbed96b0392920425b214f99df2fc1610a6056898cea7931b3fabf82e202657
% appointment-dependency: src/gpt/liturgy/roman-rite/postconciliar/roman-missal-third-edition-en-us-2011/propers/temporal/pc-s01-first-sunday-of-advent-year-b/instance/manifest.md sha256=e054344a643de91d6c50a0efe1739c000149f9eb320db8ba5b35d3045c9ca380
% appointment-dependency: src/gpt/liturgy/roman-rite/postconciliar/roman-missal-third-edition-en-us-2011/propers/temporal/pc-s01-first-sunday-of-advent-year-b/propers/verified.md sha256=79f5ae1cb162fc0a82f527d6b0ea84ae4751355264958da6c52a548025f54adc
% appointment-dependency: src/gpt/liturgy/roman-rite/postconciliar/roman-missal-third-edition-en-us-2011/propers/temporal/pc-s01-first-sunday-of-advent-year-b/research/chronology-inputs.toml sha256=10e8abeb3e5cebdeed32a483d9d595e53db4f6ca9bbe2d112a32f42c1b223e00
% appointment-dependency: src/gpt/liturgy/roman-rite/postconciliar/roman-missal-third-edition-en-us-2011/propers/temporal/shared/formularies/pc-s01-first-sunday-of-advent/propers/verified.md sha256=965cdeb11946487ad1c7979c309818f435decefc68661f1a1835f3cacd703940
% appointment-note: entrance: Cf. Psalm 25:1-3; adaptation; hebrew numbering; canonical verse queries, not appointed wording.
% appointment-note: collect: The Collect is a composed Missal text for which the formulary appoints no biblical citation.
% appointment-note: first-reading: Isaiah 63:16b-17; appointed; vulgate numbering; whole-verse query envelope, not subverse chronology or appointed wording.
% appointment-note: first-reading: Isaiah 64:1,3-8; appointed; vulgate numbering; canonical verse queries, not appointed wording.
% appointment-note: psalm: Psalm 80:2ac,3b,15-16,18-19; appointed; hebrew numbering; whole-verse query envelope, not subverse chronology or appointed wording.
% appointment-note: psalm: Psalm 80:4; response; hebrew numbering; canonical verse queries, not appointed wording.
% appointment-note: second-reading: 1 Corinthians 1:3-9; appointed; vulgate numbering; canonical verse queries, not appointed wording.
% appointment-note: acclamation: Psalm 85:8; appointed; hebrew numbering; canonical verse queries, not appointed wording.
% appointment-note: gospel: Mark 13:33-37; appointed; vulgate numbering; canonical verse queries, not appointed wording.
% appointment-note: offerings: The Prayer over the Offerings is a composed Missal text for which the formulary appoints no biblical citation.
% appointment-note: preface: The Preface I of Advent is a composed Missal text for which the formulary appoints no biblical citation.
% appointment-note: communion: Psalm 85:13; appointed; hebrew numbering; canonical verse queries, not appointed wording.
% appointment-note: after-communion: The Prayer after Communion is a composed Missal text for which the formulary appoints no biblical citation.
% appointment-note: advent-blessing: The Optional solemn blessing for Advent is a composed Missal text for which the formulary appoints no biblical citation.
% comparison-dependency: src/gpt/liturgy/roman-rite/postconciliar/roman-missal-third-edition-en-us-2011/propers/temporal/pc-s01-first-sunday-of-advent-year-b/research/chronology-profile-comparisons.toml sha256=d059aa93358b4056c2abfa79478df3df9cd68a694b67d8d6538da34376aa25df
% entrance: Entrance antiphon
\expandafter\def\csname triptychchronologyannotation@entrance\endcsname{%
  \chronologyannotationgroup{traditional-attribution}{disputed}{\textbf{Traditional attribution} -- disputed: \chronologyannotationreach{Ps.24.1}{Ps.24}{true}\chronologyannotationreach{Ps.24.2}{Ps.24}{true}\chronologyannotationreach{Ps.24.3}{Ps.24}{true}\chronologyannotationclaim{israel.monarchy.david-traditional-era}{traditional-attribution}{catholic-traditional-v1}{disputed}{reigned from 1055 to 1015 B.C.}{David (reign in the usual chronology), B.C. 1055--1015}.}%
  \space \chronologyannotationgroup{historical-setting}{research-pending}{\textbf{Historical setting} -- Occasion date unresolved.}%
  \space \chronologyannotationgroup{composition}{preferred}{\textbf{Composition}: \chronologyannotationreach{Ps.24.1}{Ps}{true}\chronologyannotationreach{Ps.24.2}{Ps}{true}\chronologyannotationreach{Ps.24.3}{Ps}{true}\chronologyannotationclaim{critical.psalms.latest-composition-boundary}{composition}{catholic-critical-v1}{preferred}{before the Maccabean period, around 165 B.C.; no individual psalm can be dated securely}{Before c. 165 B.C.}}%
}
% first-reading: First reading (U.S. Isaiah 63:16b--17, 19b; 64:2--7), corresponding Vulgate passages
\expandafter\def\csname triptychchronologyannotation@first-reading\endcsname{%
  \chronologyannotationgroup{traditional-attribution}{preferred}{\textbf{Traditional attribution}: \chronologyannotationreach{Is.63.16}{Is}{true}\chronologyannotationreach{Is.63.17}{Is}{true}\chronologyannotationreach{Is.64.1}{Is}{true}\chronologyannotationreach{Is.64.3}{Is}{true}\chronologyannotationreach{Is.64.4}{Is}{true}\chronologyannotationreach{Is.64.5}{Is}{true}\chronologyannotationreach{Is.64.6}{Is}{true}\chronologyannotationreach{Is.64.7}{Is}{true}\chronologyannotationreach{Is.64.8}{Is}{true}\chronologyannotationclaim{israel.prophets.isaias-traditional-era}{traditional-attribution}{catholic-traditional-v1}{preferred}{from the closing year of Ozias, King of Juda (740 B.C.); no prophecies appear to be later than 701 B.C.}{Isaias (ministry in Souvay's traditional account), B.C. 740--701}.}%
  \space \chronologyannotationgroup{composition}{preferred}{\textbf{Composition}: \chronologyannotationreach{Is.63.16}{Is.56 Is.57 Is.58 Is.59 Is.60 Is.61 Is.62 Is.63 Is.64 Is.65 Is.66}{true}\chronologyannotationreach{Is.63.17}{Is.56 Is.57 Is.58 Is.59 Is.60 Is.61 Is.62 Is.63 Is.64 Is.65 Is.66}{true}\chronologyannotationreach{Is.64.1}{Is.56 Is.57 Is.58 Is.59 Is.60 Is.61 Is.62 Is.63 Is.64 Is.65 Is.66}{true}\chronologyannotationreach{Is.64.3}{Is.56 Is.57 Is.58 Is.59 Is.60 Is.61 Is.62 Is.63 Is.64 Is.65 Is.66}{true}\chronologyannotationreach{Is.64.4}{Is.56 Is.57 Is.58 Is.59 Is.60 Is.61 Is.62 Is.63 Is.64 Is.65 Is.66}{true}\chronologyannotationreach{Is.64.5}{Is.56 Is.57 Is.58 Is.59 Is.60 Is.61 Is.62 Is.63 Is.64 Is.65 Is.66}{true}\chronologyannotationreach{Is.64.6}{Is.56 Is.57 Is.58 Is.59 Is.60 Is.61 Is.62 Is.63 Is.64 Is.65 Is.66}{true}\chronologyannotationreach{Is.64.7}{Is.56 Is.57 Is.58 Is.59 Is.60 Is.61 Is.62 Is.63 Is.64 Is.65 Is.66}{true}\chronologyannotationreach{Is.64.8}{Is.56 Is.57 Is.58 Is.59 Is.60 Is.61 Is.62 Is.63 Is.64 Is.65 Is.66}{true}\chronologyannotationclaim{critical.third-isaiah-oracles}{composition}{catholic-critical-v1}{preferred}{postexilic period}{After the Babylonian exile}.}%
}
% psalm: Responsorial psalm (refrain at v. 4)
\expandafter\def\csname triptychchronologyannotation@psalm\endcsname{%
  \chronologyannotationgroup{historical-setting}{research-pending}{\textbf{Historical setting} -- Occasion date unresolved.}%
  \space \chronologyannotationgroup{composition}{preferred}{\textbf{Composition}: \chronologyannotationreach{Ps.79.15}{Ps}{true}\chronologyannotationreach{Ps.79.16}{Ps}{true}\chronologyannotationreach{Ps.79.18}{Ps}{true}\chronologyannotationreach{Ps.79.19}{Ps}{true}\chronologyannotationreach{Ps.79.2}{Ps}{true}\chronologyannotationreach{Ps.79.3}{Ps}{true}\chronologyannotationreach{Ps.79.4}{Ps}{true}\chronologyannotationclaim{critical.psalms.latest-composition-boundary}{composition}{catholic-critical-v1}{preferred}{before the Maccabean period, around 165 B.C.; no individual psalm can be dated securely}{Before c. 165 B.C.}}%
}
% second-reading: Second reading (Paul writes from Ephesus to Corinth)
\expandafter\def\csname triptychchronologyannotation@second-reading\endcsname{%
  \chronologyannotationgroup{composition}{disputed}{\textbf{Composition} -- disputed: \chronologyannotationreach{1Cor.1.3}{1Cor}{true}\chronologyannotationreach{1Cor.1.4}{1Cor}{true}\chronologyannotationreach{1Cor.1.5}{1Cor}{true}\chronologyannotationreach{1Cor.1.6}{1Cor}{true}\chronologyannotationreach{1Cor.1.7}{1Cor}{true}\chronologyannotationreach{1Cor.1.8}{1Cor}{true}\chronologyannotationreach{1Cor.1.9}{1Cor}{true}\chronologyannotationclaim{composition.first-epistle-to-the-corinthians}{composition}{catholic-traditional-v1}{disputed}{(1 and 2 Corinthians; Galatians), 56}{A.D. 56}; \chronologyannotationreach{1Cor.1.3}{1Cor}{true}\chronologyannotationreach{1Cor.1.4}{1Cor}{true}\chronologyannotationreach{1Cor.1.5}{1Cor}{true}\chronologyannotationreach{1Cor.1.6}{1Cor}{true}\chronologyannotationreach{1Cor.1.7}{1Cor}{true}\chronologyannotationreach{1Cor.1.8}{1Cor}{true}\chronologyannotationreach{1Cor.1.9}{1Cor}{true}\chronologyannotationclaim{composition.first-epistle-to-the-corinthians}{composition}{catholic-traditional-v1}{disputed}{about Easter A.D. 57}{A.D. 57}.}%
}
% acclamation: Alleluia verse
\expandafter\def\csname triptychchronologyannotation@acclamation\endcsname{%
  \chronologyannotationgroup{historical-setting}{research-pending}{\textbf{Historical setting} -- Occasion date unresolved.}%
  \space \chronologyannotationgroup{composition}{preferred}{\textbf{Composition}: \chronologyannotationreach{Ps.84.8}{Ps.41 Ps.42 Ps.43 Ps.44 Ps.45 Ps.46 Ps.47 Ps.48 Ps.83 Ps.84 Ps.86 Ps.87}{true}\chronologyannotationclaim{composition.psalms-of-the-sons-of-korah}{composition}{catholic-traditional-v1}{preferred}{between the days of Isaias and the return from exile}{Between the days of Isaias and the return from exile}.}%
}
% gospel: Gospel
\expandafter\def\csname triptychchronologyannotation@gospel\endcsname{%
  \chronologyannotationgroup{narrated-event}{preferred}{\textbf{Event}: \chronologyannotationreach{Mark.13.33}{Mark.13.33-37}{false}\chronologyannotationreach{Mark.13.34}{Mark.13.33-37}{false}\chronologyannotationreach{Mark.13.35}{Mark.13.33-37}{false}\chronologyannotationreach{Mark.13.36}{Mark.13.33-37}{false}\chronologyannotationreach{Mark.13.37}{Mark.13.33-37}{false}\chronologyannotationclaim{life-of-christ.closing-charge-to-watch}{narrated-event}{catholic-traditional-v1}{preferred}{c. A.D. 29}{Jesus charges all the disciples to keep watch, c. A.D. 29 (derived)}.}%
  \space \chronologyannotationgroup{composition}{disputed}{\textbf{Composition} -- disputed: \chronologyannotationreach{Mark.13.33}{Mark}{true}\chronologyannotationreach{Mark.13.34}{Mark}{true}\chronologyannotationreach{Mark.13.35}{Mark}{true}\chronologyannotationreach{Mark.13.36}{Mark}{true}\chronologyannotationreach{Mark.13.37}{Mark}{true}\chronologyannotationclaim{composition.gospel-of-mark}{composition}{catholic-traditional-v1}{disputed}{the tenth or twelfth year after the Ascension (A.D. 38-40)}{A.D. 38--40}; \chronologyannotationreach{Mark.13.33}{Mark}{true}\chronologyannotationreach{Mark.13.34}{Mark}{true}\chronologyannotationreach{Mark.13.35}{Mark}{true}\chronologyannotationreach{Mark.13.36}{Mark}{true}\chronologyannotationreach{Mark.13.37}{Mark}{true}\chronologyannotationclaim{composition.gospel-of-mark}{composition}{catholic-traditional-v1}{disputed}{A.D. 40}{A.D. 40}; \chronologyannotationreach{Mark.13.33}{Mark}{true}\chronologyannotationreach{Mark.13.34}{Mark}{true}\chronologyannotationreach{Mark.13.35}{Mark}{true}\chronologyannotationreach{Mark.13.36}{Mark}{true}\chronologyannotationreach{Mark.13.37}{Mark}{true}\chronologyannotationclaim{composition.gospel-of-mark}{composition}{catholic-traditional-v1}{disputed}{the third year of Claudius (A.D. 43)}{A.D. 43}; \chronologyannotationreach{Mark.13.33}{Mark}{true}\chronologyannotationreach{Mark.13.34}{Mark}{true}\chronologyannotationreach{Mark.13.35}{Mark}{true}\chronologyannotationreach{Mark.13.36}{Mark}{true}\chronologyannotationreach{Mark.13.37}{Mark}{true}\chronologyannotationclaim{composition.gospel-of-mark}{composition}{catholic-traditional-v1}{disputed}{somewhere between A.D. 50 and 67}{A.D. 50--67}; \chronologyannotationreach{Mark.13.33}{Mark}{true}\chronologyannotationreach{Mark.13.34}{Mark}{true}\chronologyannotationreach{Mark.13.35}{Mark}{true}\chronologyannotationreach{Mark.13.36}{Mark}{true}\chronologyannotationreach{Mark.13.37}{Mark}{true}\chronologyannotationclaim{composition.gospel-of-mark}{composition}{catholic-traditional-v1}{disputed}{between 60 and 70}{A.D. 60--70}; \chronologyannotationreach{Mark.13.33}{Mark}{true}\chronologyannotationreach{Mark.13.34}{Mark}{true}\chronologyannotationreach{Mark.13.35}{Mark}{true}\chronologyannotationreach{Mark.13.36}{Mark}{true}\chronologyannotationreach{Mark.13.37}{Mark}{true}\chronologyannotationclaim{composition.gospel-of-mark}{composition}{catholic-traditional-v1}{disputed}{after Matthew's original-language Gospel and before Jerusalem's destruction}{After Matthew's original-language Gospel and before Jerusalem's destruction}.}%
  \space \chronologyannotationcomparisongroup{catholic-critical-v1}{composition}{composition-only}{\textbf{Composition (Catholic critical chronology, for comparison)}: \chronologyannotationreach{Mark.13.33}{Mark.1 Mark.2 Mark.3 Mark.4 Mark.5 Mark.6 Mark.7 Mark.8 Mark.9 Mark.10 Mark.11 Mark.12 Mark.13 Mark.14 Mark.15 Mark.16.1-8}{true}\chronologyannotationreach{Mark.13.34}{Mark.1 Mark.2 Mark.3 Mark.4 Mark.5 Mark.6 Mark.7 Mark.8 Mark.9 Mark.10 Mark.11 Mark.12 Mark.13 Mark.14 Mark.15 Mark.16.1-8}{true}\chronologyannotationreach{Mark.13.35}{Mark.1 Mark.2 Mark.3 Mark.4 Mark.5 Mark.6 Mark.7 Mark.8 Mark.9 Mark.10 Mark.11 Mark.12 Mark.13 Mark.14 Mark.15 Mark.16.1-8}{true}\chronologyannotationreach{Mark.13.36}{Mark.1 Mark.2 Mark.3 Mark.4 Mark.5 Mark.6 Mark.7 Mark.8 Mark.9 Mark.10 Mark.11 Mark.12 Mark.13 Mark.14 Mark.15 Mark.16.1-8}{true}\chronologyannotationreach{Mark.13.37}{Mark.1 Mark.2 Mark.3 Mark.4 Mark.5 Mark.6 Mark.7 Mark.8 Mark.9 Mark.10 Mark.11 Mark.12 Mark.13 Mark.14 Mark.15 Mark.16.1-8}{true}\chronologyannotationcomparisonclaim{critical.gospel-of-mark}{composition}{catholic-critical-v1}{catholic-critical-v1}{preferred}{artifact.united-states-conference-of-catholic-bishops.new-american-bible-revised-edition.english-usccb-web-2026-11-02.mark-introduction-web-text-272bcce7}{perhaps shortly after the year 70}{The Gospel of Mark under the NABRE's tentative modern dating proposal, perhaps shortly after A.D. 70}.}%
}
% communion: Communion antiphon
\expandafter\def\csname triptychchronologyannotation@communion\endcsname{%
  \chronologyannotationgroup{historical-setting}{research-pending}{\textbf{Historical setting} -- Occasion date unresolved.}%
  \space \chronologyannotationgroup{composition}{preferred}{\textbf{Composition}: \chronologyannotationreach{Ps.84.13}{Ps.41 Ps.42 Ps.43 Ps.44 Ps.45 Ps.46 Ps.47 Ps.48 Ps.83 Ps.84 Ps.86 Ps.87}{true}\chronologyannotationclaim{composition.psalms-of-the-sons-of-korah}{composition}{catholic-traditional-v1}{preferred}{between the days of Isaias and the return from exile}{Between the days of Isaias and the return from exile}.}%
}
% presentation: orientation-v3
\newcommand{\chronologyorientationrow}[2]{%
  \ifcsname triptychchronologyorientation@#1\endcsname
    \csname triptychchronologyorientation@#1\endcsname{#2}%
  \else
    \PackageError{triptych}{No chronology orientation row for #1}{}%
  \fi}
\expandafter\def\csname triptychchronologyorientation@row-58260a07838c47ee\endcsname#1{%
Entrance antiphon: Cf. Ps. 25:1--3 (Heb.; Vg. 24) & Traditional attribution: David's reign in the usual chronology; Hebron, then Jerusalem & c.~1055--1015 B.C. (usual chronology; Corbett also reports dates 30--50 years later)\chronologyannotationreach{Ps.24.1}{Ps.24}{true}\chronologyannotationreach{Ps.24.2}{Ps.24}{true}\chronologyannotationreach{Ps.24.3}{Ps.24}{true}\chronologyannotationclaim{israel.monarchy.david-traditional-era}{traditional-attribution}{catholic-traditional-v1}{disputed}{reigned from 1055 to 1015 B.C.}{}\\*[.1em]%
}
\expandafter\def\csname triptychchronologyorientation@row-64e8364c49076aea\endcsname#1{%
 & Historical setting & Occasion date unresolved\\*[.1em]%
}
\expandafter\def\csname triptychchronologyorientation@row-3454ae1c13c47abf\endcsname#1{%
 & Composition: Psalms & \chronologyannotationreach{Ps.24.1}{Ps}{true}\chronologyannotationreach{Ps.24.2}{Ps}{true}\chronologyannotationreach{Ps.24.3}{Ps}{true}\chronologyannotationclaim{critical.psalms.latest-composition-boundary}{composition}{catholic-critical-v1}{preferred}{before the Maccabean period, around 165 B.C.; no individual psalm can be dated securely}{Before c.~165 B.C.}\\[.3em]%
}
\expandafter\def\csname triptychchronologyorientation@row-b8de8edab8423d2d\endcsname#1{%
Responsorial psalm (refrain at v. 4): Ps. 80:2ac,3b,15--16,18--19 (Heb.; Vg. 79), Ps. 80:4 (Heb.; Vg. 79) & Historical setting & Occasion date unresolved\\*[.1em]%
}
\expandafter\def\csname triptychchronologyorientation@row-468c724c08eefcc0\endcsname#1{%
 & Composition: Psalms & \chronologyannotationreach{Ps.79.15}{Ps}{true}\chronologyannotationreach{Ps.79.16}{Ps}{true}\chronologyannotationreach{Ps.79.18}{Ps}{true}\chronologyannotationreach{Ps.79.19}{Ps}{true}\chronologyannotationreach{Ps.79.2}{Ps}{true}\chronologyannotationreach{Ps.79.3}{Ps}{true}\chronologyannotationreach{Ps.79.4}{Ps}{true}\chronologyannotationclaim{critical.psalms.latest-composition-boundary}{composition}{catholic-critical-v1}{preferred}{before the Maccabean period, around 165 B.C.; no individual psalm can be dated securely}{Before c.~165 B.C.}\\[.3em]%
}
\expandafter\def\csname triptychchronologyorientation@row-da9380ef1d93c43b\endcsname#1{%
Alleluia verse: Ps. 85:8 (Heb.; Vg. 84) & Historical setting & Occasion date unresolved\\*[.1em]%
}
\expandafter\def\csname triptychchronologyorientation@row-eeb5da3922829cb6\endcsname#1{%
 & Composition: The Psalms of the Sons of Korah & c.~740--536 B.C. (Korahite collection estimate)\chronologyannotationreach{Ps.84.8}{Ps.41 Ps.42 Ps.43 Ps.44 Ps.45 Ps.46 Ps.47 Ps.48 Ps.83 Ps.84 Ps.86 Ps.87}{true}\chronologyannotationclaim{composition.psalms-of-the-sons-of-korah}{composition}{catholic-traditional-v1}{preferred}{between the days of Isaias and the return from exile}{}\\[.3em]%
}
\expandafter\def\csname triptychchronologyorientation@row-3aa6aace7cd55134\endcsname#1{%
Communion antiphon: Ps. 85:13 (Heb.; Vg. 84) & Historical setting & Occasion date unresolved\\*[.1em]%
}
\expandafter\def\csname triptychchronologyorientation@row-d82d51fcbab70f35\endcsname#1{%
 & Composition: The Psalms of the Sons of Korah & c.~740--536 B.C. (Korahite collection estimate)\chronologyannotationreach{Ps.84.13}{Ps.41 Ps.42 Ps.43 Ps.44 Ps.45 Ps.46 Ps.47 Ps.48 Ps.83 Ps.84 Ps.86 Ps.87}{true}\chronologyannotationclaim{composition.psalms-of-the-sons-of-korah}{composition}{catholic-traditional-v1}{preferred}{between the days of Isaias and the return from exile}{}\\[.3em]%
}
\expandafter\def\csname triptychchronologyorientation@row-70fe366947726a89\endcsname#1{%
First reading (U.S. Isaiah 63:16b--17, 19b; 64:2--7), corresponding Vulgate passages: Is. 63:16b--17, Is. 64:1,3--8 & Traditional attribution: Isaias (ministry in Souvay's traditional account); Judah and Jerusalem & \chronologyannotationreach{Is.63.16}{Is}{true}\chronologyannotationreach{Is.63.17}{Is}{true}\chronologyannotationreach{Is.64.1}{Is}{true}\chronologyannotationreach{Is.64.3}{Is}{true}\chronologyannotationreach{Is.64.4}{Is}{true}\chronologyannotationreach{Is.64.5}{Is}{true}\chronologyannotationreach{Is.64.6}{Is}{true}\chronologyannotationreach{Is.64.7}{Is}{true}\chronologyannotationreach{Is.64.8}{Is}{true}\chronologyannotationclaim{israel.prophets.isaias-traditional-era}{traditional-attribution}{catholic-traditional-v1}{preferred}{from the closing year of Ozias, King of Juda (740 B.C.); no prophecies appear to be later than 701 B.C.}{B.C.~740--701}\\*[.1em]%
}
\expandafter\def\csname triptychchronologyorientation@row-10897ea45e62b321\endcsname#1{%
 & Composition: Oracles in Isaiah 56--66 & \chronologyannotationreach{Is.63.16}{Is.56 Is.57 Is.58 Is.59 Is.60 Is.61 Is.62 Is.63 Is.64 Is.65 Is.66}{true}\chronologyannotationreach{Is.63.17}{Is.56 Is.57 Is.58 Is.59 Is.60 Is.61 Is.62 Is.63 Is.64 Is.65 Is.66}{true}\chronologyannotationreach{Is.64.1}{Is.56 Is.57 Is.58 Is.59 Is.60 Is.61 Is.62 Is.63 Is.64 Is.65 Is.66}{true}\chronologyannotationreach{Is.64.3}{Is.56 Is.57 Is.58 Is.59 Is.60 Is.61 Is.62 Is.63 Is.64 Is.65 Is.66}{true}\chronologyannotationreach{Is.64.4}{Is.56 Is.57 Is.58 Is.59 Is.60 Is.61 Is.62 Is.63 Is.64 Is.65 Is.66}{true}\chronologyannotationreach{Is.64.5}{Is.56 Is.57 Is.58 Is.59 Is.60 Is.61 Is.62 Is.63 Is.64 Is.65 Is.66}{true}\chronologyannotationreach{Is.64.6}{Is.56 Is.57 Is.58 Is.59 Is.60 Is.61 Is.62 Is.63 Is.64 Is.65 Is.66}{true}\chronologyannotationreach{Is.64.7}{Is.56 Is.57 Is.58 Is.59 Is.60 Is.61 Is.62 Is.63 Is.64 Is.65 Is.66}{true}\chronologyannotationreach{Is.64.8}{Is.56 Is.57 Is.58 Is.59 Is.60 Is.61 Is.62 Is.63 Is.64 Is.65 Is.66}{true}\chronologyannotationclaim{critical.third-isaiah-oracles}{composition}{catholic-critical-v1}{preferred}{postexilic period}{After the Babylonian exile}\\[.3em]%
}
\expandafter\def\csname triptychchronologyorientation@row-5f7da3d77eeef1f3\endcsname#1{%
Gospel: Mark 13:33--37 & Narrated event: Jesus charges all the disciples to keep watch; on Mount Olivet opposite the Temple & \chronologyannotationreach{Mark.13.33}{Mark.13.33-37}{false}\chronologyannotationreach{Mark.13.34}{Mark.13.33-37}{false}\chronologyannotationreach{Mark.13.35}{Mark.13.33-37}{false}\chronologyannotationreach{Mark.13.36}{Mark.13.33-37}{false}\chronologyannotationreach{Mark.13.37}{Mark.13.33-37}{false}\chronologyannotationclaim{life-of-christ.closing-charge-to-watch}{narrated-event}{catholic-traditional-v1}{preferred}{c. A.D. 29}{c.~A.D.~29 (derived; Maas's chronology; Passion dates disputed)}\\*[.1em]%
}
\expandafter\def\csname triptychchronologyorientation@row-a64cf773c5828e9c\endcsname#1{%
 & Composition: Mark; Rome traditionally & \chronologyannotationreach{Mark.13.33}{Mark}{true}\chronologyannotationreach{Mark.13.34}{Mark}{true}\chronologyannotationreach{Mark.13.35}{Mark}{true}\chronologyannotationreach{Mark.13.36}{Mark}{true}\chronologyannotationreach{Mark.13.37}{Mark}{true}\chronologyannotationclaim{composition.gospel-of-mark}{composition}{catholic-traditional-v1}{disputed}{the tenth or twelfth year after the Ascension (A.D. 38-40)}{Proposals span A.D.~38--70}\chronologyannotationreach{Mark.13.33}{Mark}{true}\chronologyannotationreach{Mark.13.34}{Mark}{true}\chronologyannotationreach{Mark.13.35}{Mark}{true}\chronologyannotationreach{Mark.13.36}{Mark}{true}\chronologyannotationreach{Mark.13.37}{Mark}{true}\chronologyannotationclaim{composition.gospel-of-mark}{composition}{catholic-traditional-v1}{disputed}{A.D. 40}{}\chronologyannotationreach{Mark.13.33}{Mark}{true}\chronologyannotationreach{Mark.13.34}{Mark}{true}\chronologyannotationreach{Mark.13.35}{Mark}{true}\chronologyannotationreach{Mark.13.36}{Mark}{true}\chronologyannotationreach{Mark.13.37}{Mark}{true}\chronologyannotationclaim{composition.gospel-of-mark}{composition}{catholic-traditional-v1}{disputed}{the third year of Claudius (A.D. 43)}{}\chronologyannotationreach{Mark.13.33}{Mark}{true}\chronologyannotationreach{Mark.13.34}{Mark}{true}\chronologyannotationreach{Mark.13.35}{Mark}{true}\chronologyannotationreach{Mark.13.36}{Mark}{true}\chronologyannotationreach{Mark.13.37}{Mark}{true}\chronologyannotationclaim{composition.gospel-of-mark}{composition}{catholic-traditional-v1}{disputed}{somewhere between A.D. 50 and 67}{}\chronologyannotationreach{Mark.13.33}{Mark}{true}\chronologyannotationreach{Mark.13.34}{Mark}{true}\chronologyannotationreach{Mark.13.35}{Mark}{true}\chronologyannotationreach{Mark.13.36}{Mark}{true}\chronologyannotationreach{Mark.13.37}{Mark}{true}\chronologyannotationclaim{composition.gospel-of-mark}{composition}{catholic-traditional-v1}{disputed}{between 60 and 70}{}\chronologyannotationreach{Mark.13.33}{Mark}{true}\chronologyannotationreach{Mark.13.34}{Mark}{true}\chronologyannotationreach{Mark.13.35}{Mark}{true}\chronologyannotationreach{Mark.13.36}{Mark}{true}\chronologyannotationreach{Mark.13.37}{Mark}{true}\chronologyannotationclaim{composition.gospel-of-mark}{composition}{catholic-traditional-v1}{disputed}{after Matthew's original-language Gospel and before Jerusalem's destruction}{}\\*[.1em]%
}
\expandafter\def\csname triptychchronologyorientation@row-48b996b8139bea86\endcsname#1{%
 & Composition, Catholic critical comparison: Mark; Syria possible & \chronologyannotationreach{Mark.13.33}{Mark.1 Mark.2 Mark.3 Mark.4 Mark.5 Mark.6 Mark.7 Mark.8 Mark.9 Mark.10 Mark.11 Mark.12 Mark.13 Mark.14 Mark.15 Mark.16.1-8}{true}\chronologyannotationreach{Mark.13.34}{Mark.1 Mark.2 Mark.3 Mark.4 Mark.5 Mark.6 Mark.7 Mark.8 Mark.9 Mark.10 Mark.11 Mark.12 Mark.13 Mark.14 Mark.15 Mark.16.1-8}{true}\chronologyannotationreach{Mark.13.35}{Mark.1 Mark.2 Mark.3 Mark.4 Mark.5 Mark.6 Mark.7 Mark.8 Mark.9 Mark.10 Mark.11 Mark.12 Mark.13 Mark.14 Mark.15 Mark.16.1-8}{true}\chronologyannotationreach{Mark.13.36}{Mark.1 Mark.2 Mark.3 Mark.4 Mark.5 Mark.6 Mark.7 Mark.8 Mark.9 Mark.10 Mark.11 Mark.12 Mark.13 Mark.14 Mark.15 Mark.16.1-8}{true}\chronologyannotationreach{Mark.13.37}{Mark.1 Mark.2 Mark.3 Mark.4 Mark.5 Mark.6 Mark.7 Mark.8 Mark.9 Mark.10 Mark.11 Mark.12 Mark.13 Mark.14 Mark.15 Mark.16.1-8}{true}\chronologyannotationcomparisonclaim{critical.gospel-of-mark}{composition}{catholic-critical-v1}{catholic-critical-v1}{preferred}{artifact.united-states-conference-of-catholic-bishops.new-american-bible-revised-edition.english-usccb-web-2026-11-02.mark-introduction-web-text-272bcce7}{perhaps shortly after the year 70}{Perhaps shortly after A.D.~70}\\[.3em]%
}
\expandafter\def\csname triptychchronologyorientation@row-70f590fc8306682b\endcsname#1{%
Second reading (Paul writes from Ephesus to Corinth): 1 Cor. 1:3--9 & Composition: The First Epistle to the Corinthians & \chronologyannotationreach{1Cor.1.3}{1Cor}{true}\chronologyannotationreach{1Cor.1.4}{1Cor}{true}\chronologyannotationreach{1Cor.1.5}{1Cor}{true}\chronologyannotationreach{1Cor.1.6}{1Cor}{true}\chronologyannotationreach{1Cor.1.7}{1Cor}{true}\chronologyannotationreach{1Cor.1.8}{1Cor}{true}\chronologyannotationreach{1Cor.1.9}{1Cor}{true}\chronologyannotationclaim{composition.first-epistle-to-the-corinthians}{composition}{catholic-traditional-v1}{disputed}{(1 and 2 Corinthians; Galatians), 56}{Proposals span A.D.~56--57}\chronologyannotationreach{1Cor.1.3}{1Cor}{true}\chronologyannotationreach{1Cor.1.4}{1Cor}{true}\chronologyannotationreach{1Cor.1.5}{1Cor}{true}\chronologyannotationreach{1Cor.1.6}{1Cor}{true}\chronologyannotationreach{1Cor.1.7}{1Cor}{true}\chronologyannotationreach{1Cor.1.8}{1Cor}{true}\chronologyannotationreach{1Cor.1.9}{1Cor}{true}\chronologyannotationclaim{composition.first-epistle-to-the-corinthians}{composition}{catholic-traditional-v1}{disputed}{about Easter A.D. 57}{}\\[.3em]%
}

% Generated. Do not edit.
% Complete orientation-v3 table; wording, dates and order belong to the corpus.
\begin{chronologysharedtable}
\chronologyorientationrow{row-58260a07838c47ee}{}
\chronologyorientationrow{row-64e8364c49076aea}{}
\chronologyorientationrow{row-3454ae1c13c47abf}{}
\midrule
\chronologyorientationrow{row-b8de8edab8423d2d}{}
\chronologyorientationrow{row-468c724c08eefcc0}{}
\midrule
\chronologyorientationrow{row-da9380ef1d93c43b}{}
\chronologyorientationrow{row-eeb5da3922829cb6}{}
\midrule
\chronologyorientationrow{row-3aa6aace7cd55134}{}
\chronologyorientationrow{row-d82d51fcbab70f35}{}
\midrule
\chronologyorientationrow{row-70fe366947726a89}{}
\chronologyorientationrow{row-10897ea45e62b321}{}
\midrule
\chronologyorientationrow{row-4356b7782aaa8ac7}{}
\chronologyorientationrow{row-be2981581c9dd920}{}
\chronologyorientationrow{row-02b4f06cbf04508a}{}
\midrule
\chronologyorientationrow{row-70f590fc8306682b}{}
\end{chronologysharedtable}
{\footnotesize Psalm numbering: Vg. = Vulgate; Heb. = Hebrew.\par}


\subsection{Evidence for each date}
\label{sec:date-location-evidence}
{\small For each passage, as the Scriptural Date and Location table names it: what is dated, how the corpus weighs each answer, the source's own words, and the sources cited. A subject the corpus does not date is reported as such and given no date.\par}
{\small
\par\smallskip\noindent\textit{Entrance antiphon: Cf. Ps. 25:1--3 (Heb.; Vg. 24)}\par\nobreak
\begin{itemize}[leftmargin=1.2em,itemsep=1pt,topsep=1pt,parsep=0pt]
\item \textbf{Traditional attribution} (disputed). David's reign in the usual chronology: B.C.~1055--1015 (Corbett also reports writers dating the reign 30--50 years later) (``reigned from 1055 to 1015 B.C.''; The Catholic Encyclopedia, Volume 4 (New York, 1908), article 'David, King,' printed p. 642, opening paragraph). Calendar orientation: c.~1055--1015 B.C. (usual chronology; Corbett also reports dates 30--50 years later) (The Catholic Encyclopedia, Volume 4 (New York, 1908), article 'David, King,' printed p. 642, opening paragraph). Place: Hebron, then Jerusalem (The Holy Bible, Douay-Rheims (Challoner revision), Project Gutenberg eBook 1581: 2 Kings 5:5).
\item \textbf{Historical setting}. No date: the chronology corpus has not yet dated this (research pending).
\item \textbf{Composition} (preferred). Psalms: Before c.~165 B.C. (``before the Maccabean period, around 165 B.C.; no individual psalm can be dated securely''; Catholic critical chronology; New American Bible, Revised Edition --- USCCB web state (2026), Psalms introduction).
\end{itemize}
\par\smallskip\noindent\textit{Responsorial psalm (refrain at v. 4): Ps. 80:2ac,3b,15--16,18--19 (Heb.; Vg. 79), Ps. 80:4 (Heb.; Vg. 79)}\par\nobreak
\begin{itemize}[leftmargin=1.2em,itemsep=1pt,topsep=1pt,parsep=0pt]
\item \textbf{Historical setting}. As under Entrance antiphon: Cf. Ps. 25:1--3 (Heb.; Vg. 24).
\item \textbf{Composition}. As under Entrance antiphon: Cf. Ps. 25:1--3 (Heb.; Vg. 24).
\end{itemize}
\par\smallskip\noindent\textit{Alleluia verse: Ps. 85:8 (Heb.; Vg. 84)}\par\nobreak
\begin{itemize}[leftmargin=1.2em,itemsep=1pt,topsep=1pt,parsep=0pt]
\item \textbf{Historical setting}. As under Entrance antiphon: Cf. Ps. 25:1--3 (Heb.; Vg. 24).
\item \textbf{Composition} (preferred). The Psalms of the Sons of Korah: Between the days of Isaias and the return from exile (The Catholic Encyclopedia, Volume 12 (1911)). Calendar orientation: c.~740--536 B.C. (Korahite collection estimate) (The Catholic Encyclopedia, Volume 12 (1911); The Catholic Encyclopedia, Volume 8 (New York, 1910); The Catholic Encyclopedia, Volume 3 (New York, 1908)).
\end{itemize}
\par\smallskip\noindent\textit{Communion antiphon: Ps. 85:13 (Heb.; Vg. 84)}\par\nobreak
\begin{itemize}[leftmargin=1.2em,itemsep=1pt,topsep=1pt,parsep=0pt]
\item \textbf{Historical setting}. As under Entrance antiphon: Cf. Ps. 25:1--3 (Heb.; Vg. 24).
\item \textbf{Composition}. As under Alleluia verse: Ps. 85:8 (Heb.; Vg. 84).
\end{itemize}
\par\smallskip\noindent\textit{First reading (U.S. Isaiah 63:16b--17, 19b; 64:2--7), corresponding Vulgate passages: Is. 63:16b--17, Is. 64:1,3--8}\par\nobreak
\begin{itemize}[leftmargin=1.2em,itemsep=1pt,topsep=1pt,parsep=0pt]
\item \textbf{Traditional attribution} (preferred). Isaias (ministry in Souvay's traditional account): B.C.~740--701 (``from the closing year of Ozias, King of Juda (740 B.C.); no prophecies appear to be later than 701 B.C.''; The Catholic Encyclopedia, Volume 8 (New York, 1910)). Place: Judah and Jerusalem (The Holy Bible, Douay-Rheims (Challoner revision), Project Gutenberg eBook 1581: Is. 1:1).
\item \textbf{Composition} (preferred). Oracles in Isaiah 56--66: After the Babylonian exile (``postexilic period''; Catholic critical chronology; New American Bible, Revised Edition: USCCB web state (2026), Isaiah introduction, especially the final two prose paragraphs before the outline).
\end{itemize}
\par\smallskip\noindent\textit{Gospel: Mark 13:33--37}\par\nobreak
\begin{itemize}[leftmargin=1.2em,itemsep=1pt,topsep=1pt,parsep=0pt]
\item \textbf{Narrated event} (preferred). Jesus charges all the disciples to keep watch: c.~A.D.~29 (derived; Maas's chronology; Passion dates disputed) (The Holy Bible, Douay-Rheims (Challoner revision), Project Gutenberg eBook 1581: Mark 13:3, 13:4, 13:33, 13:34, 13:37; The Catholic Encyclopedia, Volume 8 (New York, 1910), article 'Chronology of the Life of Jesus Christ,' Ninth journey, heading and Holy Week paragraph: the close of the Tuesday; The Catholic Encyclopedia, Volume 8 (New York, 1910), A. J. Maas, 'Jesus Christ,' 'The Year of the Death of Christ,' printed p. 378, left column). Place: on Mount Olivet opposite the Temple (The Holy Bible, Douay-Rheims (Challoner revision), Project Gutenberg eBook 1581: Mark 13:3, 13:4, 13:37).
\item \textbf{Composition} (disputed). The Gospel of St Mark: A.D.~38--40 (``the tenth or twelfth year after the Ascension (A.D. 38-40)''; The Catholic Encyclopedia, Volume 9 (1910)); A.D.~40 (The Catholic Encyclopedia, Volume 9 (1910)); A.D.~43 (``the third year of Claudius (A.D. 43)''; The Catholic Encyclopedia, Volume 9 (1910)); A.D.~50--67 (``somewhere between A.D. 50 and 67''; The Catholic Encyclopedia, Volume 9 (1910)); A.D.~60--70 (``between 60 and 70''; The Catholic Encyclopedia, Volume 14 (New York, 1912)); After Matthew's original-language Gospel and before Jerusalem's destruction (De auctore, de tempore compositionis et de historica veritate Evangeliorum secundum Marcum et secundum Lucam (1912), responsa-i-ix). Place: Rome traditionally (The Catholic Encyclopedia, Volume 9 (1910)).
\item \textbf{Composition, Catholic critical comparison} (preferred). The Gospel of Mark under the NABRE's tentative modern dating proposal: Perhaps shortly after A.D.~70 (``perhaps shortly after the year 70''; Catholic critical chronology; New American Bible, Revised Edition: USCCB web state, 2 November 2026). Place: Syria possible (New American Bible, Revised Edition: USCCB web state, 2 November 2026).
\end{itemize}
\par\smallskip\noindent\textit{Second reading (Paul writes from Ephesus to Corinth): 1 Cor. 1:3--9}\par\nobreak
\begin{itemize}[leftmargin=1.2em,itemsep=1pt,topsep=1pt,parsep=0pt]
\item \textbf{Composition} (disputed). The First Epistle to the Corinthians: A.D.~56 (``(1 and 2 Corinthians; Galatians), 56''; The Catholic Encyclopedia, Volume 11 (New York, 1911)); A.D.~57 (``about Easter A.D. 57''; The Catholic Encyclopedia, Volume 4 (New York, 1908)).
\end{itemize}
\par}
